% ============ 33 Weil 猜想证明的完成 (原书 pp.193-196) ============
\chapter{Weil 猜想证明的完成}\label{sec:33}
\index{弱 Lefschetz 定理}
\index{Leray 谱序列}
\index{爆破}
\index{消失循环}

设 $X_0$ 是 $\mathbb{F}_q$ 上偶维数 $2m + 2 = n + 1$ 的光滑射影簇.由命题~\ref{prop:28.3}，还剩下证明：
\begin{quote}
($*$) $F$ 在 $H^{n+1}(X, \mathbb{Q}_{\ell})$ 上有理作用，且它的特征值 $\alpha$ 满足
\[
q^{n/2} < |\alpha| < q^{n/2+1}.
\]
\end{quote}
我们通过对 $m$ 作归纳证明 ($*$).

我们说有限维 $\mathbb{Q}_{\ell}$-向量空间 $V$ 的自同态 $\varphi \colon V \to V$ 满足 ($\dagger.n$)，如果它有理作用且特征值 $\alpha$ 满足
\[
q^{\frac{n-1}{2}} < |\alpha| < q^{\frac{n+1}{2}}.
\]

\begin{lemma}\label{lem:33.1}
(a) 若 $V$ 满足 ($\dagger.n$) 且 $W$ 是 $\varphi$ 稳定的 $V$ 的子空间，则 $W$ 与 $V/W$ 都满足 ($\dagger.n$).

(b) 若存在被 $\varphi$ 作用稳定的滤过
\[
V \supset V_1 \supset \cdots \supset V_r \supset 0,
\]
使得对所有的 $i$，$\varphi$ 在 $V_i/V_{i+1}$ 上定义的自同态满足 ($\dagger.n$)，则 $\varphi$ 满足 ($\dagger.n$).
\end{lemma}

\begin{proof}
读者易证的习题.
\end{proof}

与 (\ref{lem:28.1}) 证明相同的论证表明可以扩张基域 $\mathbb{F}_q$.这允许我们假定存在 $X_0$ 的在 $\mathbb{F}_q$ 上有理的 Lefschetz 铅笔；把铅笔写为 $D_0$($= \mathbb{P}^{1}$)，$S \subset D$ 为奇异超平面截的集合，$U_0 = D_0 \setminus S$.可以假定(在扩张 $\mathbb{F}_q$ 后):
\begin{enumerate}
  \item 每个 $s \in S$ 在 $\mathbb{F}_q$ 上有理，且定义 $s$ 处切锥的二次型可以(在 $\mathbb{F}_q$ 上)表示为
  \[
  Q(T_1, \ldots, T_n) = \sum_{i=1}^{m} T_i T_{i+1} + T^{2}_{2m+1};
  \]
  \item 存在 $u_0 \in \mathbb{P}^{1}(\mathbb{F}_q)$，使得纤维 $X_{u_0} \overset{\mathrm{def}}{=} \varphi^{-1}(u_0)$ 非奇异，且 $X_{u_0}$ 有定义在 $\mathbb{F}_q$ 上的非奇异超平面截 $Y_0$.
\end{enumerate}
注意 $X$ 的维数为 $n + 1 = 2m + 2$，纤维的维数为 $n = 2m + 1$，$Y_0$ 的维数为 $n - 1 = 2m$.

则沿轴 $A \cap X$ 爆破 $X$ 得到的簇 $X^{*}$ 也定义在 $\mathbb{F}_q$ 上，且有定义在 $\mathbb{F}_q$ 上的映射 $\varphi'_0 \colon X^{*}_0 \to \mathbb{P}^{1}$.我们把 $\mathbb{F}_q$ 上的射影直线写为 $\mathbb{P}^{1}$，把 $\mathbb{F}$ 上的射影直线写为 $\mathbb{P}^{1}$.

设 $u$ 是 $U$ 中映到 $u_0$($u_0$ 如 (b))的点.则 $u$ 可以看作 $U_0$ 的几何点.我们把 $\pi_1(U_0)$ 写为 $\pi^{\mathrm{tame}}_1(U_0, u)$，把 $\pi_1(U)$ 写为 $\pi^{\mathrm{tame}}_1(U, u)$.回忆有同构 $\pi_1(U_0, \bar{\xi}) \to \pi_1(U_0, u)$(良定义到共轭).

\begin{lemma}\label{lem:33.2}
只需对 $X^{*}$ 证明 ($*$).
\end{lemma}

\begin{proof}
簇 $X^{*}$ 是沿 $A \cap X$ 爆破 $X$ 得到的.因为 $A$ 横截相交 $X$，$A \cap X$ 是 $X$ 中余维数 2 的非奇异子簇，因此 $A \cap X$ 在 $X^{*}$ 中的原像 $(A \cap X)^{*}$ 同构于 $\mathbb{P}(N_{X/A \cap X})$，其中 $N_{X/A \cap X}$ 是法丛(见 Shafarevich, Vol 2, p.73).现在 (\ref{thm:23.2}) 表明
\[
H^{*}((A \cap X)^{*}, \mathbb{Q}_{\ell}) \cong H^{*}(A \cap X, \mathbb{Q}_{\ell}) \oplus H^{*-2}(A \cap X, \mathbb{Q}_{\ell})(-1).
\]
映射 $\varphi' \colon X^{*} \to X$ 紧合，因此可以对它应用紧合基变换定理(\S 17).发现典范映射 $\mathbb{Q}_{\ell} \to \varphi'_*\mathbb{Q}_{\ell}$ 是同构，$R^{2}\varphi'_*\mathbb{Q}_{\ell}$ 的支集在 $A \cap X$ 上，且对 $r \neq 0, 2$，$R^{r}\varphi'_*\mathbb{Q}_{\ell} = 0$.此外，$R^{2}\varphi'_*\mathbb{Q}_{\ell} = i_{*}(R^{2}\varphi'^{*}_0\mathbb{Q}_{\ell})$，其中 $\varphi'^{*} = \varphi'|_{(A \cap X)^{*}}$，且 $R^{2}\varphi'^{*}_0\mathbb{Q}_{\ell} = \mathbb{Q}_{\ell}(-1)$.$\varphi'$ 的 Leray 谱序列在 $E_2$ 处退化(因为 $\varphi'^{*}$ 的如此)，因此
\[
H^{*}(X^{*}, \mathbb{Q}_{\ell}) \cong H^{*}(X, \mathbb{Q}_{\ell}) \oplus H^{*-2}(A \cap X, \mathbb{Q}_{\ell})(-1).
\]
\end{proof}

由 $\varphi$ 的 Leray 谱序列，看到只需对下面三个群中的每一个证明 $F$ 的 ($\dagger.n+1$)：
\[
H^{2}(\mathbb{P}^{1}, R^{n-1}\varphi_{*}\mathbb{Q}_{\ell}), \qquad H^{1}(\mathbb{P}^{1}, R^{n}\varphi_{*}\mathbb{Q}_{\ell}), \qquad H^{0}(\mathbb{P}^{1}, R^{n+1}\varphi_{*}\mathbb{Q}_{\ell}).
\]
对 $\mathbb{P}^{1}$ 上的常值层 $V$，
\[
H^{0}(\mathbb{P}^{1}, V) = V_u;
\]
\[
H^{1}(\mathbb{P}^{1}, V) = \operatorname{Hom}(\pi_1(\mathbb{P}^{1}), V_u) = 0;
\]
\[
H^{2}(\mathbb{P}^{1}, V) = H^{0}(\mathbb{P}^{1}, \check{V}(1)^{\vee}) = V_u(-1).
\]

\textbf{群 $H^{2}(\mathbb{P}^{1}, R^{n-1}\varphi_{*}\mathbb{Q}_{\ell})$} 由 Lefschetz 铅笔的理论，$R^{n-1}\varphi_{*}\mathbb{Q}_{\ell}$ 是 $\mathbb{P}^{1}$ 上的常值层，且 $(R^{n-1}\varphi_{*}\mathbb{Q}_{\ell})_u = H^{n-1}(X_u, \mathbb{Q}_{\ell})$.因此，
\[
H^{2}(\mathbb{P}^{1}, R^{n-1}\varphi_{*}\mathbb{Q}_{\ell}) = H^{n-1}(X_u, \mathbb{Q}_{\ell})(-1).
\]
考虑
\[
0 \to j_{!}j^{*}\mathbb{Q}_{\ell} \to \mathbb{Q}_{\ell} \to i_{*}i^{*}\mathbb{Q}_{\ell} \to 0
\]
的上同调序列
\[
\cdots \to H^{n-1}_{c}(X_u \setminus Y, \mathbb{Q}_{\ell}) \to H^{n-1}(X_u, \mathbb{Q}_{\ell}) \to H^{n-1}(Y, \mathbb{Q}_{\ell}) \to \cdots.
\]
这里 $Y_0$ 是 $X_0$ 的非奇异超平面截，因此 $X_0 \setminus Y_0$ 仿射.Poincar\'e 对偶定理表明 $H^{n-1}_{c}(X_u \setminus Y, \mathbb{Q}_{\ell}) \cong H^{n+1}(X_u \setminus Y, \mathbb{Q}_{\ell})^{\vee}$，而弱 Lefschetz 定理表明它为零.因此 $H^{n-1}(X_u, \mathbb{Q}_{\ell}) \to H^{n-1}(Y, \mathbb{Q}_{\ell})$ 单射.由归纳，($\dagger$) 对 $Y$ 成立，这蕴含 $F$ 在 $H^{2}(\mathbb{P}^{1}, R^{n-1}\varphi_{*}\mathbb{Q}_{\ell})$ 上的特征值是满足
\[
q^{\frac{n}{2}} < |\alpha| < q^{\frac{n}{2}+1}
\]
的代数数 $\alpha$.

\textbf{群 $H^{0}(\mathbb{P}^{1}, R^{n+1}\varphi_{*}\mathbb{Q}_{\ell})$} 在消失循环非零的假设下，$R^{n+1}\varphi_{*}\mathbb{Q}_{\ell}$ 常值，且 $H^{0}(\mathbb{P}^{1}, R^{n+1}\varphi_{*}\mathbb{Q}_{\ell}) = H^{n+1}(X_u, \mathbb{Q}_{\ell})$.弱 Lefschetz 定理表明 Gysin 映射
\[
H^{n-1}(Y, \mathbb{Q}_{\ell})(-1) \to H^{n+1}(X_u, \mathbb{Q}_{\ell})
\]
满射，可以再对 $Y$ 应用归纳假设.

\textbf{群 $H^{1}(\mathbb{P}^{1}, R^{n}\varphi_{*}\mathbb{Q}_{\ell})$} 最后，必须处理 $H^{1}(\mathbb{P}^{1}, R^{n}\varphi_{*}\mathbb{Q}_{\ell})$.设 $V = (R^{n}\varphi_{*}\mathbb{Q}_{\ell})_u$.我们有滤过
\[
V \supset E \supset E \cap E^{\perp} \supset 0.
\]
由 Picard-Lefschetz 公式推出 $\pi_1(U)$ 在 $V/E$ 与 $E \cap E^{\perp}$ 上平凡作用.特别地，这是 $\pi_1(U)$-模的滤过，从而对应 $U$ 上层 的滤过：
\[
V \supset E \supset E \cap E^{\perp} \supset 0
\]
应用 $j_{*}$，得到滤过
\[
R^{n}\varphi_{*}\mathbb{Q}_{\ell} \supset j_{*}E \supset j_{*}(E \cap E^{\perp}) \supset 0.
\]
\[
\text{常值} \qquad \text{有趣的} \qquad \text{常值}
\]
我们用了 $R^{n}\varphi_{*}\mathbb{Q}_{\ell} \xrightarrow{\sim} j_{*}j^{*}R^{n}\varphi_{*}\mathbb{Q}_{\ell}$.

\textbf{商 $E/E \cap E^{\perp}$} 我们希望把主引理(定理~\ref{thm:30.6})应用于 $E/E \cap E^{\perp}$.

\begin{lemma}\label{lem:33.3}[有理性引理]
对所有闭点 $x \in U_0$，$F_x$ 在 $E/E \cap E^{\perp}$ 上的作用有理.
\end{lemma}

\begin{proof}
我们把证明推迟到下一小节.(证明相当复杂；见 Deligne 1974, \S 6.)
\end{proof}

杯积配对
\[
H^{n}(X_u, \mathbb{Q}_{\ell}) \times H^{n}(X_u, \mathbb{Q}_{\ell}) \to H^{2n}(X_u, \mathbb{Q}_{\ell}) \cong \mathbb{Q}_{\ell}(-n)
\]
斜对称(因为 $n$ 奇数)且非退化.因为它是典范的(事实上，它来自 $R^{n}\varphi_{*}\mathbb{Q}_{\ell}$ 上的典范配对)，它在 $\pi_1(U_0)$ 下不变.回忆 $V = H^{n}(X_u, \mathbb{Q}_{\ell})$，因此可以把配对改写为
\[
V \times V \to \mathbb{Q}_{\ell}(-n).
\]
现在 $E$ 是 $V$ 的 $\pi_1(U_0)$-不变子空间，$E^{\perp}$ 是 $E$ 对这个配对的 正交补.因此，配对诱导非退化配对
\[
\psi \colon E/E \cap E^{\perp} \times E/E \cap E^{\perp} \to \mathbb{Q}_{\ell}(-n).
\]
最后，Kazhdan 与 Margulis 的定理~\ref{thm:32.6} 表明 $\pi_1(U)$ 的像在 $\operatorname{Sp}(E/E \cap E^{\perp}, \psi)$ 中开.于是，主引理~\ref{thm:30.6} 表明 $F$ 在 $j_{*}(E/E \cap E^{\perp})$ 上的作用满足 ($\dagger.n+1$).

\textbf{常值商} 回忆 $E/E \cap E^{\perp}$ 是简单 $\pi_1$-模.因此，要么它非零(且没有消失循环在 $E \cap E^{\perp}$ 中)，要么它为零且 $E \subset E^{\perp}$.

第一种情形：没有消失循环在 $E \cap E^{\perp}$ 中.这意味着对任何 $s \in S$，$E^{I_s}$ 在 $E$ 中有余维数 1.因为 $V^{I_s}$ 在 $V$ 中有余维数 1，且 $I_s$ 在 $V/E$ 上平凡作用，序列
\[
0 \to E \to V \to V/E \to 0
\]
在取 $I_s$ 不变量后仍正合.这蕴含当我们对 $U$ 上局部常值层的相应序列
\[
0 \to E \to V \to V/E \to 0
\]
应用 $j_{*}$ 时，在 $\mathbb{P}^{1}$ 上得到的序列
\[
0 \to j_{*}E \to R^{n}\varphi_{*}\mathbb{Q}_{\ell} \to j_{*}(V/E) \to 0
\]
正合.因为 $j_{*}(V/E)$ 常值，映射
\[
H^{1}(\mathbb{P}^{1}, j_{*}E) \to H^{1}(\mathbb{P}^{1}, R^{n}\varphi_{*}\mathbb{Q}_{\ell})
\]
满射，因此只需对 $H^{1}(\mathbb{P}^{1}, j_{*}E)$ 证明 ($\dagger.n+1$).

同样，因为 $E^{I_s}$ 在 $E$ 中有余维数 1，序列
\[
0 \to E \cap E^{\perp} \to E \to E/E \cap E^{\perp} \to 0
\]
在取 $I_s$ 不变量后仍正合.因此，
\[
0 \to j_{*}(E \cap E^{\perp}) \to j_{*}E \to j_{*}(E/E \cap E^{\perp}) \to 0
\]
正合，且因为 $j_{*}(E \cap E^{\perp})$ 常值，
\[
H^{1}(\mathbb{P}^{1}, j_{*}E) \to H^{1}(\mathbb{P}^{1}, j_{*}(E/E \cap E^{\perp}))
\]
单射，从而 $H^{1}(\mathbb{P}^{1}, j_{*}E)$ 的 ($\dagger.n+1$) 由 $H^{1}(\mathbb{P}^{1}, j_{*}(E/E \cap E^{\perp}))$ 的 ($\dagger.n+1$) 推出.

第二种情形：消失循环在 $E \cap E^{\perp}$ 中，即 $E \subset E^{\perp}$.在这种情况下，定义 $\mathcal{F}$ 为商 $R^{n}\varphi_{*}\mathbb{Q}_{\ell}/j_{*}E^{\perp}$.则我们有两个正合序列
\[
0 \to j_{*}E^{\perp} \to R^{n}\varphi_{*}\mathbb{Q}_{\ell} \to \mathcal{F} \to 0
\]
\[
0 \to \mathcal{F} \to j_{*}j^{*}\mathcal{F} \to \bigoplus_{s} \mathbb{Q}_{\ell}(m - n)_s \to 0.
\]
层 $j_{*}E^{\perp}$ 与 $j_{*}j^{*}\mathcal{F}$ 常值，因此相应的上同调序列是
\[
0 \to H^{1}(\mathbb{P}^{1}, R^{n}\varphi_{*}\mathbb{Q}_{\ell}) \to H^{1}(\mathbb{P}^{1}, \mathcal{F})
\]
\[
\bigoplus_{s \in S} \mathbb{Q}_{\ell}(m - n) \to H^{1}(\mathbb{P}^{1}, \mathcal{F}) \to 0.
\]
因为 $F$ 在 $\mathbb{Q}_{\ell}(m - n)$ 上按 $q^{n-m} = q^{\frac{n+1}{2}}$ 作用.

这完成了 Weil 猜想的证明.

\subsection*{有理性引理的证明}

暂略.

\subsection*{消失循环退化时的证明}

由于消失循环(到符号)共轭---见定理~\ref{thm:32.3}---它们要么全为零，要么全非零.上面我们假定它们非零(见 \S 32).在消失循环全为零的情形，证明更容易(但暂略).
